% Standalone document; compile twice with pdflatex.
% The preparation invariant is inherited from the 120-state construction.
% The state-37 invariant below has a direct proof from the transition table.
\documentclass[11pt]{article}
\usepackage[T1]{fontenc}
\usepackage{lmodern}
\usepackage[margin=1in]{geometry}
\usepackage{amsmath,amssymb,amsthm,booktabs}
\usepackage[hidelinks]{hyperref}
\usepackage[nameinlink,capitalise]{cleveref}
\newtheorem{theorem}{Theorem}[section]
\newtheorem{lemma}[theorem]{Lemma}
\newtheorem{corollary}[theorem]{Corollary}
\crefname{theorem}{Theorem}{Theorems}
\crefname{lemma}{Lemma}{Lemmas}
\crefname{corollary}{Corollary}{Corollaries}

\title{A Further Reduction from 120 to 119 States}
\author{Joseph M. Shunia}
\date{September 28, 2026}
\hypersetup{
  pdftitle={A Further Reduction from 120 to 119 States},
  pdfauthor={Joseph M. Shunia}
}

\begin{document}
\maketitle

\section{Proof of the reduction}
\label{sec:119reduction}

Let $M_{120}$ denote the canonical 120-state machine, with working states
$Q=\{0,\ldots,119\}$, initial state $0$, and external halt target $H$.
Write
\[
  \delta(q,a)=(b,D,r)
\]
when state $q$, reading $a$, writes $b$, moves in direction
$D\in\{L,R\}$, and enters $r$. All state numbers in the following
reachability argument refer to $M_{120}$.
For a configuration $(q,h,\tau)$, the tape is
$\tau\colon\mathbb Z\to\{0,1\}$ and the head is at $h$.
The preparation-state argument below invokes the established bootstrap
and register-backend results of the original 120-state construction.

\begin{lemma}[Preparation-state restriction]
\label{lem:119prep}
Every visit to state $6$ in the blank-tape computation of $M_{120}$
scans zero.
\end{lemma}

\begin{proof}
State $6$ is the canonical image of the backend state
\texttt{1b.reg.prep\_1}. The finite bootstrap certificate
in the original manuscript (the lemma titled ``Bootstrap certificate'')
reaches the first source decision without halting.
It therefore cannot execute the completed transition
$\delta(6,1)=(1,R,H)$ during initialization.

After initialization, the register normal form and the primitive-operation
argument of the original manuscript's ``Primitive register semantics''
lemma apply. On entry to
\texttt{1b.reg.prep\_1}, the head is one cell to the right of the newly
written preparation marker, in the separating zero gap. Hence the scanned
cell is zero. Each primitive operation restores normal form; the control
simulation and restart protocol preserve this assertion at all subsequent
operations, by that manuscript's ``Backend simulation'' theorem,
``Natural restart of \texttt{main}'' lemma, and ``Dispatcher
specialization'' proposition.
The all-tapes state quotient and the subsequent 121-to-120 projection
preserve the tape and head position along the blank-tape execution, so
the same restriction holds in $M_{120}$.
This uses the established backend invariant, with no assumption about
whether the source computation eventually halts.
\end{proof}

\begin{lemma}[Return-state restriction]
\label{lem:119return}
Every visit to state $37$ in the blank-tape computation of $M_{120}$
scans one.
\end{lemma}

\begin{proof}
We give a direct argument from the transition table. Set
\[
\begin{split}
 A&=\{39,48,74,80,90,93,94,113,117\},\\
 B&=\{72,75,102,112\}.
\end{split}
\]
For each state $q\in A\cup B$, there is exactly one incoming transition.
Its source $p$ is listed below; in every case the instruction is
$\delta(p,1)=(1,R,q)$.
\begin{center}
\begin{tabular}{rr@{\qquad}rr}
\toprule
Target $q$ & Source $p$ & Target $q$ & Source $p$\\
\midrule
39 & 33 & 72 & 59\\
48 & 40 & 75 & 61\\
74 & 60 & 102 & 95\\
80 & 65 & 112 & 108\\
90 & 76 & & \\
93 & 81 & & \\
94 & 82 & & \\
113 & 109 & & \\
117 & 114 & & \\
\bottomrule
\end{tabular}
\end{center}
Since none of these states is the initial state, every configuration in
one of them follows such a transition. The cell immediately to the left
of the new head is the cell just written. Consequently
\begin{equation}
\label{eq:119leftone}
 q\in A\cup B\quad\Longrightarrow\quad\tau(h-1)=1.
\end{equation}

The complete lists of incoming transitions to states $86$ and $89$ are
\[
\begin{aligned}
 \delta(72,1)&=(1,R,86),\\
 \delta(75,1)=\delta(102,1)=\delta(112,1)&=(1,R,89).
\end{aligned}
\]
Their source states belong to $B$. Thus \eqref{eq:119leftone}, together
with the newly written one, gives
\begin{equation}
\label{eq:119twoleft}
 q\in\{86,89\}\quad\Longrightarrow\quad
 \tau(h-2)=\tau(h-1)=1.
\end{equation}

The complete list of incoming transitions to state $32$ is
\[
\begin{aligned}
 \delta(29,1)=\delta(115,1)&=(1,R,32),\\
 \delta(86,1)=\delta(89,1)&=(0,L,32).
\end{aligned}
\]
The right-moving transitions leave a newly written one immediately to
the left of the head. For a left-moving transition, the cell immediately
to the left of the new head was two cells to the left of the old head,
and is one by \eqref{eq:119twoleft}. Therefore
\begin{equation}
\label{eq:119state32}
 q=32\quad\Longrightarrow\quad\tau(h-1)=1.
\end{equation}

Finally, the incoming transitions to state $37$ are exactly
\[
 \delta(p,1)=(0,L,37)\quad(p\in A\cup\{32\}),
 \qquad \delta(37,0)=(0,L,37).
\]
Every entry from a different state therefore lands on a one, by
\eqref{eq:119leftone} and \eqref{eq:119state32}.
If a configuration in state $37$ scanning zero existed, consider the
first such configuration. It is not initial and cannot arise from a
different state. Its predecessor must therefore use the transition
$(37,0)$, contradicting the choice of the first occurrence.
Thus $(37,0)$ never occurs.
\end{proof}

\begin{theorem}[Reduction to 119 working states]
\label{thm:119reduction}
There is a deterministic binary Turing machine $M_{119}$ with 119
working states whose blank-tape run agrees with that of $M_{120}$ at
every step in tape contents and head position, with control states
related by a fixed projection. The two machines halt at exactly the
same time, or both run forever.
\end{theorem}

\begin{proof}
Define the state projection by $\pi(H)=H$ and
\[
 \pi(q)=
 \begin{cases}
 q,&0\le q\le36,\\
 6,&q=37,\\
 q-1,&38\le q\le119.
 \end{cases}
\]
Its image on working states is $\{0,\ldots,118\}$.
For every $q\notin\{6,37\}$ and $a\in\{0,1\}$, define
\[
 \delta'(\pi(q),a)=(b,D,\pi(r))
 \quad\text{when}\quad\delta(q,a)=(b,D,r).
\]
At the merged state, define
\[
 \boxed{\delta'(6,0)=(0,R,8),\qquad
        \delta'(6,1)=(0,L,10).}
\]
These are respectively the original instructions
$\delta(6,0)$ and $\delta(37,1)$, with their targets projected.
The construction gives one instruction for each of the
$119\cdot2=238$ state/symbol pairs.

By \cref{lem:119prep,lem:119return}, every state/read pair occurring
in the original blank-tape run belongs to
\[
 D=\bigl(Q\times\{0,1\}\bigr)
       \setminus\{(6,1),(37,0)\}.
\]
For every $(q,a)\in D$, the construction satisfies
\begin{equation}
\label{eq:119projection}
 \delta(q,a)=(b,D_0,r)
 \quad\Longrightarrow\quad
 \delta'(\pi(q),a)=(b,D_0,\pi(r)).
\end{equation}
There are exactly 238 such equations.

Induct on execution time. Initially the tapes are blank, the head
positions agree, and $\pi(0)=0$. Suppose the runs currently have
identical tapes and head positions, with states $q$ and $\pi(q)$.
They read the same symbol $a$, and the original run uses a pair in $D$.
Equation \eqref{eq:119projection} gives identical writes and head
moves, and projected successor states. Hence the same relation holds
after the next step. Since $\pi(H)=H$ and no working state maps to $H$,
entry to the halt target occurs simultaneously.
\end{proof}

\begin{corollary}
\label{cor:119rh}
The machine $M_{119}$ halts on the blank tape if and only if the
Riemann hypothesis is false. It has exactly one halting transition,
namely $(76,1)$.
\end{corollary}

\begin{proof}
The halting equivalence follows from \cref{thm:119reduction} and the
established RH equivalence of $M_{120}$.
The two halting transitions in the original table are $(6,1)$ and
$(77,1)$. The first is replaced at the merged state; the second is
retained with state number $\pi(77)=76$. No other transition is
changed into a halting transition.
\end{proof}

The reduction is valid for the blank-tape computation. Its two
reachability restrictions, and therefore the reduction itself, do not
assume the Riemann hypothesis. In particular, preservation holds even
if the original machine eventually reaches its arithmetic-failure halt.

\end{document}
